\section{Model Genealogy}
\label{sec:model-genealogy}

Level 1.2 and Level 1.2b are not a single design decision; they sit at the
end of an eight-milestone development sequence in which each mechanism was
introduced, defaulted to an exact no-op, and validated before the next
layer was added. Level 1.3 and Level 1.3b, the historical extension this
paper's second half reports, continue that same discipline directly:
every historical mechanism they add --- pressure-transmission channels,
entrenchment, dual legitimacy ledgers --- was itself introduced, defaulted
to a verified no-op, and mechanism-validated before any discovery or causal
experiment consumed it. This section summarizes the full ten-milestone
sequence to make clear what was already established, and under what
conditions, before each discovery pipeline reported later in this paper was
ever run. Full detail for every milestone below is in the source
repository's own per-level scientific notes and ODD addenda (see
\cref{sec:reproducibility}); this section reports only what is load-bearing
for interpreting the results that follow.

\begin{table}[htbp]
\centering
\caption{Model layers and mechanisms by milestone. Every mechanism
defaults to a verified, byte-for-byte no-op unless explicitly activated.}
\label{tab:mechanism-layer}
\adjustbox{max width=\textwidth}{\begin{tabular}{@{}lll@{}}
\toprule
\textbf{Milestone} & \textbf{Layer / mechanism added} & \textbf{Default state} \\
\midrule
Level 0 / 0.1 & Bilateral cooperation/defection, punishment, imitation, exit; resource pool & always active \\
Level 0.2     & Viability state machine, fitness modes, participation accounting & no-op (\texttt{viability\_mode="none"}) \\
Level 0.3 / 0.3b & Independent ecological policy; local patches; audited regime classifier & no-op (\texttt{ecology\_topology="global\_pool"}) \\
Level 1.0     & Proto-institution: propose/adopt/monitor/sanction/amend/dissolve & no-op (\texttt{rule\_formation\_enabled=False}) \\
Level 1.1     & Delegation, capacity, discretion, accountability, contributions, coordination & no-op (\texttt{delegation\_enabled=False}, etc.) \\
Level 1.1b    & Mechanism-substream RNG; audited delegation classifier & no-op (\texttt{rng\_mode="legacy\_single\_stream"}) \\
Level 1.2     & Local observation, transmission, migration, scaling; phenotype vector \& blind discovery & no-op (5 mechanism flags default off) \\
Level 1.2b    & Ecological ledger, feasibility ratios, feature/recurrence audits, stratified rediscovery & no new fields; reuses \texttt{level12\_complete} \\
Level 1.3     & Pressure-transmission channels (security/bargaining/coordination), entrenchment, dual legitimacy ledgers, $S_{\text{obs}}$/$Z_{\text{obs}}$ composites & no-op (channel flags default off) \\
Level 1.3b    & Sham-controlled history-erasure intervention, purpose-designed discovery/confirmation, measurement-weighting audit & no new fields; validation milestone only \\
\bottomrule
\end{tabular}
}
\end{table}

\subsection{Level 0: atomic relational mechanisms}
\label{sec:level0}

The base model (Level 0 and its Level~0.1 ecological-closure extension) is
a heterogeneous population interacting over a dynamic network via a
repeated cooperation/defection payoff grammar, with imperfect observation,
memory and forgiveness, optional costly punishment, strategy imitation
with low-rate mutation, and costly relationship exit and formation. Level
0.1 closes the resource loop with a shared, bounded environment pool
governed by logistic regeneration and mandatory subsistence costs, and
makes relationship dynamics genuinely ongoing rather than an
initialization-time artifact. Every stochastic draw in the codebase passes
through one explicit, seeded random-number generator, making every run
exactly reproducible from its configuration.

\subsection{Ecological closure and viability}
\label{sec:ecological-closure}

Level 0.2 added full participation-opportunity accounting, a disableable
agent viability state machine (distress, dormancy, and, only if
explicitly enabled, removal and replacement), configurable imitation
fitness signals, explicit ecological supply/demand balance accounting, and
operational (never "attractor"-labeled) dynamical classification
definitions. This milestone's own investigation into Level 0.1's pool
collapse and near-universal measured cooperation found both candidate
explanations it tested --- selective-participation bias and
relative-fitness invariance to universal impoverishment --- genuinely
present, with regime-dependent relative weight.

\subsection{Endogenous ecological behavior}
\label{sec:endogenous-ecology}

Level 0.3 gave agents a second, structurally independent behavioral
dimension --- an ecological action (extraction, maintenance, or restraint)
governed by its own policy, never derived from the social
cooperate/defect decision --- and an opt-in spatially disaggregated
alternative to the single global pool: fixed-membership local patches.
Four causal-intervention counterfactuals (no environmental effect, and
frozen, randomized, or shuffled ecological policy) were built so that any
observed persistence could be attributed to actual agent behavior, not
merely to the bare environmental equation. Level 0.3b then re-audited Level
0.3's own regime classifier with a stricter, burn-in-aware, multi-criterion
test and a prospectively registered confirmation pipeline; the audited
classifier found none of the original "oscillatory" or "intermittent"
labels survived across 144 point-seed runs, and established, via a
paired, prespecified causal comparison, that endogenous adaptive
ecological behavior causally worsens ecological persistence relative to
frozen or randomized counterfactuals. This causal result --- not a
discretionary preference --- was the two-route decision gate that licensed
proceeding to institutional mechanisms at all.

\subsection{Proto-institutions, delegation, accountability, and coordination}
\label{sec:proto-institutions}

Level 1.0 introduced the proto-institution substrate directly in response
to Level 0.3b's demonstrated coordination failure: a locally shared,
patch-attached rule that can be proposed, voluntarily adopted, imperfectly
monitored and enforced, amended, dissolved, and --- possibly ---
reconstructed, entirely from local information. Level 1.0's own
counterfactual comparison found persistent proto-institutions emerged
unforced at every tested ecological candidate, alongside a
directionally consistent (not, at that milestone's sample size,
statistically confirmed) pattern in which real monitoring and
sanctioning infrastructure was associated with \textit{lower} final
ecological occupancy than no institution at all. Level 1.1 then built the
generative bridge from a single static rule to delegation: administrative
capacity accumulated only from realized task performance (never a static
trait), discretion and its corresponding appropriation opportunity,
accountability (audit, recall, appeal), opt-in collective contributions,
and minimal bilateral interpatch coordination, deliberately independent of
delegation. Level 1.1b then validated these mechanisms directly ---
artifact-independent re-verification, a per-mechanism random-number
substream architecture for causal ablations, and an audited, more
conservative delegation classifier reported alongside the original for
direct comparison, never replacing it.

\subsection{Why Level 1.2 was the appropriate discovery gate}
\label{sec:why-level12}

Only once local institutional formation, delegation, accountability, and
coordination were each independently built, defaulted to exact no-ops, and
mechanism-validated in isolation did the project introduce the mechanisms
that make institutions capable of \textit{propagating and competing for
persistence across space} --- local observation, component-wise
transmission, migration and membership selection, differential
persistence, and institutional scaling --- and, simultaneously, the
observer-only apparatus needed to ask, without bias toward any answer,
whether recurrent institutional configurations exist at all: a
ten-dimensional phenotype vector, temporal feature extraction, a blind
discovery pipeline, and a twelve-criterion recurrence test. This is Level
1.2. Running an outcome-blind discovery search before this point would have
been premature (institutions could not yet propagate or compete); running
it very much later, without first establishing that every prior mechanism
defaults to a verified no-op and is individually validated, would have
made any discovered (or undiscovered) structure impossible to attribute to
specific, understood mechanisms. Level 1.2 is where those two
requirements are jointly first satisfied.

\subsection{Why Level 1.3 was the appropriate historical-extension gate}
\label{sec:why-level13}

Level 1.2b's own closing synthesis (\cref{sec:level12b-verdict}) established
that the synchronic substrate is generatively sufficient but not
phenotype-convergent, and explicitly motivated, without implementing or
testing, a historically conditioned extension
(\cref{eq:historical-decomposition-intro}). Only once that synchronic
result was independently replicated on fresh data, and only once the
project's own ecological ledger, feasibility measures, and recurrence
machinery were validated at the 576-run scale, did the project introduce
history as a generative mechanism: three pressure-transmission channels
(security, bargaining, coordination-complexity), each wired through two of
six candidate transmission primitives shared across channels
(\texttt{pressure\_adjusted\_probability}, \texttt{pressure\_adjusted\_%
scale}); institutional entrenchment (age-dependent resistance to
amendment/dissolution); and dual performance/procedural legitimacy ledgers,
each independently costly and fallible. Every one of these mechanisms
defaults to the identical Level 1.2b behavior when disabled, verified by
regression test, so that Level 1.3's own results can be attributed to the
specific historical mechanisms added, not to an undisclosed change in the
substrate Level 1.2b already validated. Running a historical-pressure
search before Level 1.2b's own synchronic null was independently replicated
would have made it impossible to tell whether any resulting structure came
from history or from unresolved ecological confounding; Level 1.3 is where
those two requirements are jointly first satisfied for the historical
extension, exactly as Level 1.2 was for the synchronic one.

\subsection{Level 1.3b: purpose-designed validation and causal intervention}
\label{sec:why-level13b}

Level 1.3's own discovery evidence (\cref{sec:level13}) was pooled and
post hoc: every headline finding rested on the same 216- or 228-run pool
assembled across Stages 3-5/7, with only one genuinely independent
fresh-data test (Stage 7's 12-point confirmation), which is exactly the one
claim that failed to reproduce. Level 1.3b closes this gap not by
re-analyzing Level 1.3's own pool but by running a purpose-designed,
preregistered, eight-hypothesis program with its own discovery pool
(Stage 4, $n=240$), its own disjoint fresh-confirmation pool (Stage 8,
$n=96$), and, for the first time in this project, a sham-controlled causal
intervention stage (Stage 6) capable of directly testing whether erasing or
altering a run's own realized history changes its trajectory, validated
against a byte-identical control that itself introduces zero spurious
divergence. \Cref{sec:level13b} reports this program in full.
